\scp{The Data}{01-03}
This book gleans data from many published
and unpublished sources.
The increasing availability of online resources,
including Dutch-era literature,
has made such gleaning possible to a breadth
and depth that was not possible even a decade ago.
This study also includes previously unpublished data from our own fieldnotes,
and from other scholars who have generously shared their own work with us.
These are mentioned in the Acknowledgments
and the Bibliography.

Appendix \ref{ch:LanLis} lists the sources from which we have gleaned
data from at least 517 languages/speech varieties.
425 are Austronesian
and 91 are from various non-Austronesian (Papuan) families.
292 of these varieties are found within our
primary target region of Linguistic Wallacea,
including 34 Papuan \tsc{Timor-Alor-Pantar} languages,
but excluding SHWNG,
\tsc{North Halmahera} languages
and other Papuan languages of the Bird’s Head
and Cenderawasih Bay region.
This volume has 77 Figures
and 334 Tables comparing different types of data from many languages.
These tallies do not count 195 numbered examples/lists
providing different types of data from many
languages or the numerous in-text examples found throughout the volume.
There are also 332 footnotes
and 444 table notes.

Our intent is that this study is data-rich.
We try to present the data in ways that make it easy for others to see the patterns that we are seeing.
And to lay it out so that the fuller data,
including exceptions, is presented for alternate analyses.

To our knowledge,
nobody has previously attempted a study of this breadth
and depth of the historical phonology of the Austronesian languages in our target region.
On the other hand,
a study of this breadth
and depth has only recently become a possibility.